% ===========================================================================
\section{Ideals}\label{sec:ideals}
\warmth{0}\quad\whisper{The algebra side. An ideal is the algebraic shadow of a geometric shape.}

\begin{strand}
The objects of nonlinear algebra live in the \keyword{polynomial ring}
$K[\xb]=K[x_1,\dots,x_n]$ over a field $K$ (favorite: $K=\Q$; also $\R,\C,\overline K,
\mathbb F_q$). A polynomial is a finite $K$-combination of monomials $\xb^{\mathbf a}
=x_1^{a_1}\cdots x_n^{a_n}$, and its \keyword{degree} is $\max\{|\mathbf a|:c_{\mathbf a}\ne0\}$.
The hero of the book is one cubic surface; meet it now, it returns everywhere.
\end{strand}

\begin{example}[the running cubic, $n=3$]\label{ex:cayley}
\[
  f \;=\; \det\!\begin{pmatrix}1&x&y\\ x&1&z\\ y&z&1\end{pmatrix}
    \;=\; 2xyz-x^2-y^2-z^2+1 .
\]
$\V(f)\subset\R^3$ is a cubic surface (Cayley's nodal cubic). Its four \keyword{singular points}
are the common zeros of $f$ and its three partials $\partial f/\partial x=2yz-2x$, etc.
\end{example}
\begin{yourturn}
Solve $2yz-2x=2xz-2y=2xy-2z=0$ together with $f=0$. Check the four singular points are
$(1,1,1),(1,-1,-1),(-1,1,-1),(-1,-1,1)$. \recall{Why must a singular point also satisfy $f=0$ here? (Euler's identity for homogeneous-ish $f$ -- or just check.)}
\workspace[3]
\end{yourturn}

\block{What an ideal is}

\begin{definition}[ideal, Def.\ 1.1]\label{def:ideal}
All rings here are commutative with $1$. A nonempty subset $I\subseteq R$ is an \keyword{ideal} if
\begin{enumerate}[label=(\alph*),leftmargin=2em,itemsep=0pt]
  \item $f\in R$ and $g\in I$ \;$\Rightarrow$\; $\fillin[2.5cm]\in I$ \quad(absorbs multiplication);
  \item $f,g\in I$ \;$\Rightarrow$\; $\fillin[2cm]\in I$ \quad(closed under addition).
\end{enumerate}
Equivalently, $I=\ker\phi$ for some ring homomorphism $\phi:R\to S$.
\end{definition}

\recall{Why is $\{$even integers$\}$ an ideal of $\Z$? Which homomorphism has it as kernel?}%
Given $\mathcal F\subseteq R$, the smallest ideal containing it is $\gen{\mathcal F}$, the ideal
\keyword{generated} by $\mathcal F$ (all polynomial-coefficient combinations). The quotient $R/I$
is a ring (Prop.\ 1.2), built from congruence classes $f+I$.

\begin{proposition}[four operations, Prop.\ 1.3]\label{prop:ops}
If $I,J\subseteq R$ are ideals, so are the sum $I+J$, the intersection $I\cap J$, the product
$IJ=\gen{fg:f\in I,g\in J}$, and the colon (quotient) $(I:J)=\{f\in R:fJ\subseteq I\}$.
\end{proposition}
\begin{proof}
We check $(I:J)$ is an ideal. For (a): $f\in R,\ g\in(I:J)\Rightarrow (fg)J=f(gJ)\subseteq fI\subseteq I$.
For (b): $f,g\in(I:J)\Rightarrow (f+g)J\subseteq fJ+gJ\subseteq I+I=I$. \TODO{do the same two checks for $I\cap J$}
\end{proof}

\begin{yourturn}
($n=1$, Example 1.4.) In $\Q[x]$ let $I=\gen{(x+2)^3}$ and $J=\gen{(x-1)(x+2)}$. Compute, as
principal ideals (factor first!):
\[
  \begin{aligned}
    I\cap J&=\fillin[2.6cm], & IJ&=\fillin[2.6cm],\\[2pt]
    I+J&=\fillin[2.6cm],     & (I:J)&=\fillin[2.6cm].
  \end{aligned}
\]
Moral: \emph{arithmetic in $\Q[x]$ is just like arithmetic in $\Z$} ($\cap=\mathrm{lcm}$, $+=\gcd$).
\end{yourturn}

\block{Four kinds of ideal, read off the quotient}

\begin{definition}[the property dictionary]\label{def:primes}
For an ideal $I\subseteq R$ and its quotient $R/I$:
\end{definition}
\noindent\begin{tabularx}{\linewidth}{@{}l l L@{}}
\toprule
{\color{madder}$I$ is\dots} & {\color{indigo}definition} & \multicolumn{1}{l}{\color{indigo}\dots iff $R/I$ is}\\
\midrule
maximal  & no proper ideal contains $I$           & a \fillin[2cm]\\
prime    & $fg\in I\Rightarrow f\in I$ or $g\in I$ & an \fillin[3cm]\\
radical  & $f^s\in I\Rightarrow f\in I$            & free of \fillin[2.5cm] elements\\
primary  & $fg\in I,\ g\notin I\Rightarrow f^s\in I$ & one where zero divisors are nilpotent\\
\bottomrule
\end{tabularx}

\begin{proposition}[implications, Prop.\ 1.6]\label{prop:impl}
$I$ maximal $\Rightarrow$ prime $\Rightarrow$ radical, and prime $\Rightarrow$ primary.
None reverses; radical $+$ primary $=$ prime; any intersection of primes is radical.
\end{proposition}
\begin{proof}
Maximal $\Rightarrow$ prime since a field has no zero divisors. Prime $\Rightarrow$ radical:
take $g=f^{s-1}$ and induct on $s$. \TODO{prove radical $+$ primary $\Rightarrow$ prime (assume $fg\in I$, $f\notin I$; use primary then radical)}
\end{proof}

\recall{Match the non-reversing counterexamples in $\R[x,y]$: $\gen{x^2}$, $\gen{x(x-1)}$, $\gen{x}$ to ``primary not radical'', ``radical not primary'', ``prime not maximal''.}

\begin{strand}
Back to the cubic (Example~\ref{ex:cayley}). Its singular locus is $\Ideal=\gen{yz-x,\,xz-y,\,xy-z}$,
a \emph{radical} ideal that factors as an intersection of \keyword{five} maximal ideals; by the
Chinese Remainder Theorem $\R[x,y,z]/\Ideal\cong\R^5$. The number five is the five complex points
of $\V(\Ideal)$, four of them the real nodes. Algebra (an intersection of maximal ideals) has just
\emph{counted geometry} (a finite set of points). That is the dictionary firing for the first time.
\end{strand}

\loose{Exercise 18 (book): identify the maximal, prime, radical, primary ideals of $R=\Z$. Same four words, older ring.}
